%\documentclass[a4paper, oneside, openany]{book}
\documentclass[a4paper]{article}

%**************************************************************
% Importazione package
%************************************************************** 

% modifica i margini
%\usepackage[top=3.1cm, bottom=3.1cm, left=2.2cm, right=2.2cm]{geometry}

% specifica con quale codifica bisogna leggere i file
\usepackage[utf8]{inputenc}

% per scrivere in italiano e in inglese;
% l'ultima lingua (l'italiano) risulta predefinita
\usepackage[english, italian]{babel}

% imposta lo stile italiano per i paragrafi
\usepackage{parskip}

% numera anche i paragrafi
\setcounter{secnumdepth}{4}

% elenca anche i paragrafi nell'indice
\setcounter{tocdepth}{4}

% permetti di definire dei colori
\usepackage[usenames,dvipsnames]{color}

% permette l'inserimento di url e di altri tipi di collegamento
\usepackage[colorlinks=true]{hyperref}

\hypersetup{
    colorlinks=true, % false: boxed links; true: colored links
    citecolor=black,
    filecolor=black,
    linkcolor=black, % color of internal links
    urlcolor=Maroon  % color of external links
}

% permette al comando \url{...} di andare a capo a metà di un link
\usepackage{breakurl}

% immagini
\usepackage{graphicx}

% permette di riferirsi all'ultima pagina con \pageref{LastPage}
\usepackage{lastpage}

% tabelle su più pagine
\usepackage{longtable}

% tabelle con il campo X per riempire lo spazio rimanente sulla riga
\usepackage{tabularx}

% personalizza l'intestazione e piè di pagina
\usepackage{fancyhdr}

% permette di inserire caratteri speciali
\usepackage{textcomp}

\fancypagestyle{plain}{
	% cancella tutti i campi di intestazione e piè di pagina
	\fancyhf{}

	\lhead{\GroupName{} \texttwelveudash \ProjectName{}}
	\chead{}
	\rhead{\slshape \leftmark}

	\lfoot{\DocTitle \\ v\DocVersion}
	\rfoot{\thepage\ di \pageref{LastPage}}

	% Visualizza una linea orizzontale in cima e in fondo alla pagina
	\renewcommand{\headrulewidth}{0.3pt}
	\renewcommand{\footrulewidth}{0.3pt}
}
\pagestyle{plain}

% Per inserire del codice sorgente formattato
\usepackage{listings}

\lstset{
  extendedchars=true,          % lets you use non-ASCII characters
  inputencoding=utf8,   % converte i caratteri utf8 in latin1, richiede \usepackage{listingsutf8} anzichè listings
  basicstyle=\ttfamily,        % the size of the fonts that are used for the code
  breakatwhitespace=false,     % sets if automatic breaks should only happen at whitespace
  breaklines=true,             % sets automatic line breaking
  captionpos=t,                % sets the caption-position to top
  commentstyle=\color{mygreen},   % comment style
  frame=none,               % adds a frame around the code
  keepspaces=true,            % keeps spaces in text, useful for keeping indentation of code (possibly needs columns=flexible)
  keywordstyle=\bfseries,     % keyword style
  numbers=none,               % where to put the line-numbers; possible values are (none, left, right)
  numbersep=5pt,              % how far the line-numbers are from the code
  numberstyle=\color{mygray}, % the style that is used for the line-numbers
  rulecolor=\color{black},    % if not set, the frame-color may be changed on line-breaks within not-black text (e.g. comments (green here))
  showspaces=false,           % show spaces everywhere adding particular underscores; it overrides 'showstringspaces'
  showstringspaces=false,     % underline spaces within strings only
  showtabs=false,             % show tabs within strings adding particular underscores
  stepnumber=5,               % the step between two line-numbers. If it's 1, each line will be numbered
  stringstyle=\color{red},    % string literal style
  tabsize=4,                  % sets default tabsize
  firstnumber=1      % visualizza i numeri dalla prima linea
}

% Permetti di utilizzare il grassetto per i caratteri Typewriter (per es. il font di \code{...} e \file{...})
\usepackage[T1]{fontenc}
\usepackage{lmodern}
%%%%%%%%%%%%%%
%  COSTANTI  %
%%%%%%%%%%%%%%

% In questa prima parte vanno definite le 'costanti' utilizzate da due o più documenti.

\newcommand{\GroupName}{\emph{SteakHolders}}
\newcommand{\ProjectName}{\emph{MaaP}}

\newcommand{\Proponente}{CoffeeStrap}
\newcommand{\Committente}{Prof. Tullio Vardanega}

% TODO I seguenti 4 sono deprecati, sono da togliere.
% Usare i 4 definiti sopra, con la maiuscola all'inizio!
\newcommand{\groupName}{\GroupName{}- FIXME: DEPRECATO}
\newcommand{\projectName}{\ProjectName{}- FIXME: DEPRECATO}
\newcommand{\proponente}{\Proponente{}- FIXME: DEPRECATO}
\newcommand{\committente}{\Committente{}- FIXME: DEPRECATO}

% La versione dei documenti deve essere definita qui in global, perchè serve anche agli altri documenti
\newcommand{\VersioneG}{??.??.??}
\newcommand{\VersionePQ}{??.??.??}
\newcommand{\VersioneNP}{1.1.5}
\newcommand{\VersioneAR}{1.1.1}

% Quando serve riferirsi a ``Nome del Documento + ultima versione x.y.z'' usiamo queste costanti:
\newcommand{\Glossario}{\emph{Glossario v\VersioneG{}}}
\newcommand{\PianoDiQualifica}{\emph{Piano di Qualifica v\VersionePQ{}}}
\newcommand{\NormeDiProgetto}{\emph{Norme di Progetto v\VersioneNP}}

\newcommand{\ScopoDelProdotto}{
	Lo scopo del progetto è la realizzazione di un \glossario{framework} per generare interfacce web di amministrazione dei dati di \glossario{business} basato su \glossario{stack} \glossario{Node.js} e \glossario{MondoDB}. L'obbiettivo è quello di semplificare il processo di implementazione di tali interfacce che lo sviluppatore, appoggiandosi alla produttività del framework MaaP, potrà generare in maniera semplice e veloce ottenendo quindi un considerevole risparmio di tempo e di sforzo. Il fruitore finale delle pagine generate sarà infine l'esperto di business che potrà visualizzare, gestire e modificare le varie entità e dati residenti in \glossario{MongoDB}.
	Il prodotto atteso si chiama \glossario{MaaP} ossia \emph{MongoDB as an admin Platform}.
}

%%%%%%%%%%%%%%
%  FUNZIONI  %
%%%%%%%%%%%%%%

% In questa seconda parte vanno definite le 'funzioni' utilizzate da due o più documenti.

% Serve a dare la giusta formattazione alle parole presenti nel glossario
% il nome del comando \glossary è già usato da LaTeX
\newcommand{\glossario}[1]{\mbox{#1\ped{\ped{G}}}}

% Serve a dare la giusta formattazione al codice inline
\newcommand{\code}[1]{\texttt{#1}}

% Serve a dare la giusta formattazione a tutte le path presenti nei documenti
\newcommand{\file}[1]{\texttt{#1}}

% Permette di andare a capo all'interno di una cella in una tabella
\newcommand{\multiLineCell}[2][c]{\begin{tabular}[#1]{@{}l@{}}#2\end{tabular}}

% Genera automaticamente la pagina di copertina
\newcommand{\makeFrontPage}{
  \begin{titlepage}
  \begin{center}

  \begin{center}
  \includegraphics[width=10cm]{../../modello/steakman.png}
  \end{center}
  
  \vspace{33pt}

  \begin{Huge}
  \textbf{\DocTitle{}}
  \end{Huge}
  
  \textbf{\emph{Gruppo} \GroupName{} \, \texttwelveudash \, \emph{Progetto} \ProjectName{}}
  
  \vspace{11pt}

  \bgroup
  \def\arraystretch{2}
  \begin{tabular}{ r|l }
    \multicolumn{2}{c}{\textbf{Informazioni sul documento}} \\
    \hline
    \textbf{Versione} & \DocVersion{} \\
    \textbf{Redazione} & \multiLineCell[t]{\DocRedazione{}} \\
    \textbf{Verifica} & \multiLineCell[t]{\DocVerifica{}} \\
    \textbf{Approvazione} & \multiLineCell[t]{\DocApprovazione{}} \\
    \textbf{Uso} & \DocUso{} \\
    \textbf{Distribuzione} & \multiLineCell[t]{\DocDistribuzione{}} \\
  \end{tabular}
  \egroup

  \vspace{22pt}

  \textbf{Descrizione} \\
  \DocDescription{}

  \end{center}
  \end{titlepage}
}

%%%%%%%%%%%%%%
%  COSTANTI  %
%%%%%%%%%%%%%%

% In questa prima parte vanno definite le 'costanti' utilizzate soltanto da questo documento.
% Devono iniziare con una lettera maiuscola per distinguersi dalle funzioni.

\newcommand{\DocTitle}{Norme di Progetto}
\newcommand{\DocVersion}{\VersioneNP{}}
\newcommand{\DocRedazione}{
	Federico Poli \\
	Luca De Franceschi \\
	Giacomo Fornari \\
	Nicolò Tresoldi
}
\newcommand{\DocVerifica}{Serena Girardi}
\newcommand{\DocApprovazione}{Nicolò Tresoldi}
\newcommand{\DocUso}{Interno}
\newcommand{\DocDistribuzione}{
	\Committente{} \\
	Gruppo \GroupName{}
}

% La descrizione del documento
\newcommand{\DocDescription}{
 Questo documento descrive le regole, gli strumenti e le procedure adottate dal gruppo \GroupName{} per la realizzazione del progetto \ProjectName{}.
}

%%%%%%%%%%%%%%
%  FUNZIONI  %
%%%%%%%%%%%%%%

% In questa seconda parte vanno definite le 'funzioni' utilizzate soltanto da questo documento.


\begin{document}

\makeFrontPage

\section*{Registro delle modifiche}

\small{
\begin{tabularx}{\textwidth}{|c|c|l|X|}
 \hline \textbf{Versione} & \textbf{Data} & \textbf{Persone coinvolte} & \textbf{Descrizione} \\

 % IN ORDINE DALLA MODIFICA PIÙ RECENTE ALLA PIÙ VECCHIA
 % TODO Compilarlo in modo verosimile riferendosi alle date del Piano di Progetto
 
 \hline 1.1.5 & 2013-12-05 & Nome Cognome &
 Stesura dello scheletro iniziale e del capitolo ``Introduzione''.\\

 \hline 1.1.4 & 2013-12-06 & Nome Cognome &
 ...\\

 \hline 1.1.3 & 2013-12-08 & Nome Cognome &
 ...\\

 \hline 1.1.2 & 2013-12-09 & Nome Cognome &
 ...\\

 \hline 1.1.1 & 2013-12-10 & Nome Cognome &
 ...\\

 \hline
\end{tabularx}
}


\clearpage
\tableofcontents

\letteraGlossario{A}

\definizione{Abstract syntax tree}
Rappresentazione ad albero della struttura sintattica astratta del codice sorgente. Ogni nodo dell'albero denota un costrutto nel codice.

\definizione{Accidentali}
Problematiche non intrinseche alla produzione, ma che ne sono direttamente collegate.

\definizione{ActiveAdmin}
\glossario{Ruby on Rails} plugin per la generazione di interfacce di amministrazione dati in modo rapido.

\definizione{Agile}
Modello di ciclo di sviluppo software nato alla fine degli anni '90, si basa su quattro principi:
\begin{itemize} 
 \item L'eccessiva rigidità ostacola l'emergere del valore;
 \item Concentrarsi più sul software che sulla documentazione;
 \item La collaborazione con gli \glossario{stakeholder} è fondamentale;
 \item La capacità di adattamento al cambiare delle situazioni è importante.
\end{itemize} 
Si fa riferimento al loro manifesto \url{http://agilemanifesto.org/}

\definizione{Amazon}
Società di commercio elettronico statunitense che offre tra le altre cose una piattaforma di cloud computing.

\definizione{Amazon AWS}
Amazon Web Server è un servizio di \glossario{Amazon} che mette a disposizione agli utenti un web \glossario{server} il cui piano tariffario prevede 750 ore mensili gratuite per il primo anno dopo l'iscrizione.

\definizione{Angular.js}
Angular.js è un \glossario{framework} \glossario{open source} \glossario{Javascript}, mantenuto da Google, utilizzato per creare componenti \glossario{front-end}.

\definizione{API}
Application programming interface, sono un insieme di procedure che un sistema software rende accessibile a terzi per interfacciarsi ad esso.

\definizione{Application logic}
Parte del sistema che si occupa di processare i dati, e che non interagisce con l'\glossario{utente}. Solitamente risiede nel \glossario{back-end}, ma può anche risiedere nel \glossario{front-end}.

\definizione{Applicazione MaaP}
Maap inteso come applicativo disponibile all'utente.

\definizione{Aspell}
Software \glossario{open source} per il controllo ortografico.

\definizione{Autoformazione}
Si tratta di ore spese dai componenti del gruppo per apprendere le tecnologie utilizzate, tali ore non sono a carico del \glossario{proponente}.

\letteraGlossario{B}

\definizione{Back-end}
Parte di un software che elabora i dati generati dal front end e non interagisce con l'\glossario{utente}. Spesso si trova in una macchina diversa da quella del front-end.

\definizione{Baseline}
Uno stato dell'insieme di documenti e del software in cui ogni elemento è verificato ed approvato.

\definizione{Base di dati}
Archivio di dati in cui le informazioni contenute sono strutturate seguendo una logica, può essere relazionale o non-relazionale.

\definizione{Best practice}
Insieme di tecniche, pattern e paradigmi che formano il miglior modo per procedere.

\definizione{Black box}
Contestualizzato nell ambito del software testing si riferisce a quei test che valutano solo input e output, è un approccio agnostico rispetto al funzionamento interno del componente testato.

\definizione{Boilerplate code}
Sezioni di codice ridonandante.

\definizione{Bottom-up}
Il bottom up prende corpo dal punto di partenza (bottom) ovvero dalla situazione iniziale; considera l'obiettivo finale, induce a costruire un percorso sequenziale organizzato in passaggi successivi in cui l'ancoraggio tra traguardi intermedi e obiettivo finale è generalmente ricercato in modo intuitivo (euristico).

\definizione{Branch}
Istruzione che quando eseguita potrebbe portare all'inizio di una diversa sequenza di istruzioni.

\definizione{Broken windows theory}
Teoria sociale per cui non vengono tollerate le piccole trasgressioni che, se trascurate, potrebbero generare fenomeni di emulazione. 

\definizione{Bug}
Errore nel codice di un software.

\definizione{Burndown}
Rappresentazione grafica del lavoro da fare su un progetto nel tempo. Di solito il lavoro rimanente è indicato sull'asse verticale e il tempo sull'asse orizzontale.

\definizione{Business}
Inteso come dominio di Business, l'insieme di tutti i dati che riguardano un dato campo.

\definizione{By correction}
Ottenere la correttezza di un software procedendo per correzioni, ovvero applicando un metodo iterativo. Si tratta di un metodo errato in quanto fa perdere molto tempo e non garantisce a priori la correttezza finale.

\letteraGlossario{C}

\definizione{C++}
Linguaggio di programmazione orientato agli oggetti, con tipizzazione statica, sviluppato nel 1983 come un miglioramento del linguaggio C.

\definizione{Callback}
Funzione, o blocco di codice che viene passata come parametro ad un'altra funzione.

\definizione{CamelCase}
Metodo per scrivere parole composte o frasi unendo tutte le parole tra loro, ma lasciando le loro iniziali maiuscole.

\definizione{Chain of Responsibility}
\glossario{Design pattern} comportamentale che fa in modo di non accoppiare il mittente di una richiesta con il ricevitore, dando a più oggetti la possibilità di gestire la richiesta.

\definizione{Chrome}
Browser basato su WebKit sviluppato da Google. La versione più recente alla stesura di questo documento è la 31.x

\definizione{Client-Server}
Architettura di rete nella quale un computer client si connette ad un server per richiedere un certo servizio, quale ad esempio la condivisione di una certa risorsa, appoggiandosi ad una suite di protocolli.

\definizione{Collection}
In \glossario{MongoDB} è un insieme di \glossario{documents}.

\definizione{Collection-Index}
Insieme di pagine generate da \ProjectName{} contenenti la visualizzazione in forma tabellare dell'elenco di tutti i documenti della \glossario{Collection} \glossario{MongoDB}.

\definizione{Collection-Show}
Insieme di pagine generate da \ProjectName{} contenenti tutte le coppie chiavi-valore.

\definizione{Committente}
Il committente è la figura che commissiona un lavoro, in questo caso il progetto \ProjectName{}.

\definizione{Consuntivo}
Rendiconto  dei risultati di un dato periodo di attività.

\definizione{Cookie}
I cookie HTTP sono righe di testo usate per eseguire autenticazioni automatiche, tracciatura di sessioni e memorizzazione di informazioni specifiche riguardanti gli utenti che accedono al server. Nel dettaglio, sono stringhe di testo di piccola dimensione inviate da un server ad un Web client e poi rimandati indietro dal client al server ogni volta che il client accede alla stessa porzione dello stesso dominio web.

\definizione{CRUD}
Acronimo per Create, Retrieve, Update, Delete. Sono le funzioni standard di una \glossario{base di dati}.

\definizione{CSV}
Acronimo per Comma-Separated Values, è un formato di file basato su file di testo utilizzato per l'importazione ed esportazione di una tabella di dati.

\letteraGlossario{D}

\definizione{Dashboard}
Pagina web che rappresenta lo stato corrente di un'applicazione con un interfaccia semplice ed immediata.

\definizione{Database}
Vedi Base di dati.

\definizione{Dati di business}
Vedi Business.

\definizione{Debugging}
Attività volta ad eliminare i \glossario{bug} interni ad un sistema. 

\definizione{Deployment}
Traducibile con spiegamento, si intende la consegna al cliente del software.

\definizione{Dependency Injection}
\`E un design pattern della Programmazione orientata agli oggetti il cui scopo è quello di semplificare lo sviluppo e migliorare la testabilità di software di grandi dimensioni.

\definizione{Design pattern}
Soluzione progettuale generale per la risoluzione di un problema ricorrente.

\definizione{Documents}
\glossario{MongoDB} è un \glossario{database} a documenti. I dati al suo interno sono dunque inseriti in documenti, che possono essere paragonati alle tabelle nel mondo dei \glossario{database} relazionali, hanno però una struttura meno rigida e sono codificati in BSON.

\definizione{Document}
Equivale ad una riga di una tabella. Vedi la definizione di Documents.

\definizione{Documento}
File binario o testuale, può contenere vari tipi di informazione come testo, immagini, ecc.

\definizione{DSL}
Domain Specific Language è un linguaggio di programmazione o un linguaggio di specifica dedicato a particolari problemi di un dominio o a una particolare tecnica di rappresentazione.

\definizione{Duck Typing}
Stile di tipizzazione dinamica dove la semantica di un oggetto è determinata dall'insieme corrente dei suoi metodi e delle sue proprietà anziché dal fatto di estendere una particolare classe o implementare una specifica interfaccia.

\letteraGlossario{E}

\definizione{Ember.js}
Ember.js è un \glossario{framework} \glossario{open source} \glossario{Javascript}, basato sul pattern architetturale MVC, viene utilizzato dal gruppo per creare le componenti \glossario{front-end} di \ProjectName{}.

\definizione{Ereditarietà}
L'ereditarietà è uno dei concetti fondamentali nel paradigma di programmazione a oggetti. Consiste in una relazione che il linguaggio di programmazione, o il programmatore stesso, stabilisce tra due classi.

\definizione{Esperti di business}
Persone esperte del dominio di \glossario{business}.

\definizione{Express}
\glossario{Framework} per \glossario{Node.js}.

\definizione{Event-driven}
Paradigma di programmazione, nei programmi scritti utilizzando la tecnica a eventi il flusso del programma è largamente determinato dal verificarsi di eventi esterni. 

\letteraGlossario{F}

\definizione{Facade}
Nella programmazione ad oggetti indica un oggetto che permette, attraverso un'interfaccia più semplice, l'accesso a sottosistemi che espongono interfacce complesse e molto diverse tra loro, nonché a blocchi di codice complessi.

\definizione{Factory}
\glossario{Design pattern} creazionale che fornisce un interfaccia per creare una famiglia di oggetti senza specificare classi concrete.

\definizione{Failure}
Traducibile con malfunzionamento. In un software si ha una failure quando il comportamento del sistema discosta da quanto era stato progettato.

\definizione{Facilitatore}
È qualcuno che aiuta un gruppo di persone a capire i loro obiettivi comuni e che le assiste nel raggiungerli senza prendere posizione nella discussione.

\definizione{Firefox}
Browser \glossario{open source} prodotto da Mozilla Foundation.

\definizione{Footer}
Traducibile con piè di pagina è una sezione di pagina posizionata sotto il testo vero e proprio

\definizione{Framework}
Struttura di supporto su cui un applicativo può essere progettato.
Un framework comprende librerie di codice, convenzioni di sviluppo e una serie di strumenti di supporto allo sviluppo.

\definizione{Framework MaaP}
MaaP inteso come Framework, vedi la relativa definizione.

\definizione{Front-end}
Parte di un sistema software che gestisce l'interazione con l'utente o con sistemi esterni che producono dati di ingresso.

\definizione{FunGoStudios}
\glossario{StartUp} che si occupa di mobile games.

\definizione{Funzionale}
Vedi \glossario{linguaggi funzionali}.

\letteraGlossario{G}

\definizione{Gantt}
Diagramma di Gantt, è uno strumento di supporto alla gestione dei progetti, prende il nome da Henry Laurence Gantt.
Permette di distribuire le attività pianificate nell'arco temporale stabilito.

\definizione{GanttProject}
\glossario{Tool} online che permette la creazione di diagrammi di \glossario{Gantt} in modo cooperativo.

\definizione{Garbage}
In informatica denota porzioni di memoria che non dovranno più essere successivamente utilizzate dalle applicazioni e che quindi occupano inutilmente spazio.

\definizione{GitHub}
Servizio web di hosting per lo sviluppo di progetti software che usa il sistema di controllo di versione Git.

\definizione{Google Calendar}
Servizio gratuito offerto da Google che permette di creare e condividere calendari.

\definizione{GUI}
Graphical User Interface, è un tipo di interfaccia che consente all'utente di interagire con la macchina manipolando oggetti grafici convenzionali.

\definizione{Gulpease}
L'indice Gulpease è un indice di leggibilità di un testo tarato sulla lingua italiana. Rispetto ad altri ha il vantaggio di utilizzare la lunghezza delle parole in lettere anziché in sillabe, semplificandone il calcolo automatico.

\definizione{Gzip}
Gzip è un programma libero per la compressione dei dati. Il suo nome è la contrazione di GNU zip. Fu inizialmente creato da Jean-Loup Gailly e Mark Adler. La versione 0.1 fu rilasciata pubblicamente il 31 ottobre 1992.

\letteraGlossario{H}

\definizione{Halstead}
Metriche software introdotte da Maurice Howard Halstead nel 1977.

\definizione{Heroku}
Offre una piattaforma in cloud per ospitare applicazioni, supportando svariati linguaggi.

\definizione{HTML/CSS/JS}
Linguaggi comunemente utilizzati nell'ambito web.

\definizione{HTTP}
Acronimo di HyperText Transfer Protocol è un protocollo utilizzato come sistema per la trasmissione di informazioni sul web.
Le specifiche del protocollo sono gestite dal World Wide Web Consortium (W3C).

\letteraGlossario{I}

\definizione{IDE}
Integrated development environment,  è un software che, in periodo di programmazione, aiuta i programmatori nello sviluppo del codice sorgente di un programma.

\definizione{Index-page}
Vedi Collection Index.

\definizione{Issue}
Sinonimo di \glossario{Ticket}.

\definizione{I/O}
Input/Output.

\definizione{Incapsulamento}
Nell'ambito della programmazione, si definisce incapsulamento (o encapsulation) la tecnica di nascondere il funzionamento interno - deciso in fase di progetto - di una parte di un programma, in modo da proteggere le altre parti del programma dai cambiamenti che si produrrebbero in esse nel caso che questo funzionamento fosse difettoso, oppure si decidesse di implementarlo in modo diverso. Per avere una protezione completa è necessario disporre di una robusta interfaccia che protegga il resto del programma dalla modifica delle funzionalità soggette a più frequenti cambiamenti.

\letteraGlossario{J}

\definizione{Java}
Linguaggio di programmazione orientato agli oggetti, creato da Sun Microsystems.

\definizione{Javascript}
Linguaggio di scripting orientato agli oggetti comunemente usato nella programmazione Web.

\definizione{Jenkins}
Strumento \glossario{open source} per l'integrazione continua, scritto in linguaggio \glossario{Java}. 

\definizione{JSON}
Acronimo di \glossario{JavaScript} Object Notation, è un formato adatto per lo scambio dei dati in applicazioni \glossario{client-server} basato sul linguaggio \glossario{JavaScript}.

\letteraGlossario{L}

\definizione{LAN}
Local Area Network, è una rete informatica di collegamento tra più computer, che copre un'area limitata, come un'abitazione, una scuola o un'azienda.

\definizione{Linea di comando}
Tipologia di interfaccia utente caratterizzata da un'interazione di tipo testuale tra utente ed elaboratore: l'utente impartisce comandi testuali in input mediante tastiera e riceve risposte testuali in output dall'elaboratore mediante display.

\definizione{Linguaggi funzionali}
Paradigma di programmazione in cui il flusso di esecuzione del programma assume la forma di una serie di valutazioni di funzioni matematiche. Il punto di forza principale di questo paradigma è la mancanza di effetti collaterali delle funzioni, il che comporta una più facile \glossario{verifica} della correttezza e della mancanza di \glossario{bug} del programma e la possibilità di una maggiore ottimizzazione dello stesso.

\definizione{Link simbolico}
Particolare tipo di file contenente un percorso relativo od assoluto ad un file o directory a cui fa riferimento.

\definizione{Log}
File in cui si tengono registrate le attività compiute per esempio da un'applicazione, da un \glossario{server}, o da un interprete di comandi.

\definizione{Logger}
In informatica è un componente non intrusivo usato nei periodi di test di un prodotto software con lo scopo di registrare dei dati sull'esecuzione per favorire l'analisi dei risultati.

\letteraGlossario{M}

\definizione{MaaP}
\glossario{Framework} per generare interfacce web di amministrazione dei \glossario{dati di business} basato su stack \glossario{Node.js} e \glossario{MongoDB}.

\definizione{MaaS}
\ProjectName{} come servizio web.

\definizione{Map-Reduce}
Si tratta di un modello di programmazione per processare grandi quantità di dati in modo parallelo e distribuito su cluster di calcolo.

\definizione{Makefile}
File in cui vengono elencati una serie di istruzioni che verranno eseguiti tramite il comando \code{make}.

\definizione{Markdown}
Linguaggi per la formattazione del testo. Vengono implementati da software utilizzati dal gruppo come \glossario{GitHub} e \glossario{TeamworkPM}.

\definizione{Markup}
Linguaggio basato su un meccanismo sintattico che consente di definire e controllare il significato degli elementi contenuti in un documento o in un testo.

\definizione{Microsoft Project}
Software di pianificazione sviluppato e venduto da Microsoft. È  uno strumento per assistere i project manager nella pianificazione, nell'assegnazione delle risorse, nella verifica del rispetto dei tempi, nella gestione dei budget e nell'analisi dei carichi di lavoro.

\definizione{Middleware}
Insieme di programmi che fungono da intermediari tra diverse applicazioni e componenti software. Sono spesso utilizzati come supporto per sistemi distribuiti complessi.

\definizione{Milestone}
Termine inglese che letteralmente significa pietra miliare. Viene utilizzato nella pianificazione e gestione di progetti per indicare il raggiungimento di obiettivi stabiliti in periodo di definizione del progetto. 
Le milestone indicano quindi importanti traguardi intermedi nello svolgimento del progetto.

\definizione{Moduli}
\glossario{Unità} di lavoro di programmazione affidabile ad una singola persona. Tali moduli dovrebbero avere dimensione predefinita per evitare carichi di lavoro troppo grandi. Si può lavorare a più moduli in parallelo.

\definizione{MOM}
Acronimo di Message-oriented middleware, è un'infrastruttura \glossario{client-server} che, distribuendo un'applicazione tra più piattaforme eterogenee, ne incrementa l'interoperabilità, la portabilità e la flessibilità, semplificando lo sviluppo di applicazioni che usano sistemi operativi e protocolli di rete diversi, permettendo al programmatore di ignorare i dettagli degli stessi.
Il software MOM risiede sia sul client che sul server e consente l'esecuzione di chiamate asincrone fra gli stessi.

\definizione{MongoDB}
Sistema gestionale di basi di dati non relazionale, orientato ai documenti, di tipo \glossario{NoSQL}. Il linguaggio utilizzato per la gestione dei dati è \glossario{JavaScript}, del quale sfrutta in particolare la notazione BSON.

\definizione{Mongoose}
Utilizzato per creare una struttura logica nei documenti di \glossario{MongoDB}.

\definizione{MVC}
Acronimo di Model View Controller, è un \glossario{Design Pattern} architetturale che divide l'applicazione in tre parti interconnesse, separando l'interfaccia utente, la logica e i dati.

\definizione{MVW}
Acronimo di Model View Whatever  \glossario{Design Pattern} architetturale derivato da \glossario{MVC} ed utilizzato dal \glossario{framework} \glossario{Angular.js}.
Il significato del termine Whatever dichiarato dagli \glossario{sviluppatori} è "Whatever works for you".

\definizione{MVVM}
Acronimo di Model View View-Controller, è un \glossario{Design Pattern} derivato da \glossario{MVC}, in cui la view non comunica direttamente con il model.

\letteraGlossario{N}

\definizione{Namespace}
Collezione di nomi di entità, definite dal programmatore, omogeneamente usate in uno o più file sorgente.  Lo scopo dei namespace è quello di evitare confusione ed equivoci nel caso siano necessarie molte entità con nomi simili, fornendo il modo di raggruppare i nomi per categorie. 

\definizione{ngrok}
Software che permette di interfacciare un web service locale a internet in modo sicuro, su una data porta.

\definizione{Node.js}
Piattaforma software utilizzata per creare applicazioni distribuite facilmente scalabili.
Node.js utilizza \glossario{JavaScript} come linguaggio di scripting e gestisce le attese I/O in modo asincrono.

\definizione{NoSQL}
NoSQL è un movimento che promuove sistemi software dove la persistenza dei dati caratterizzata dal fatto di non utilizzare il modello relazionale, tipicamente usato dai \glossario{database} tradizionali. L'espressione NoSQL fa riferimento al linguaggio SQL, che è il più comune linguaggio di interrogazione dei dati nei \glossario{database} relazionali, qui preso a simbolo dell'intero paradigma relazionale.
Questi archivi di dati tipicamente non richiedono uno \glossario{schema} fisso, evitano spesso le operazioni di unione e puntano a scalare orizzontalmente.

\definizione{Node Package Manager}
Node Package Manager è il package manager ufficiale per Node.js, gestisce le dipendenze e permette di installare le applicazioni registrate nel sistema Node Package Manager. Viene eseguito tramite il comando \code{npm}.

\letteraGlossario{O}

\definizione{Open Source}
Software i cui autori ne permettono e favoriscono il libero studio e l'apporto di modifiche da parte di altri programmatori indipendenti. Questo è realizzato mediante l'applicazione di apposite licenze d'uso.

\definizione{OOP}
La programmazione orientata agli oggetti (OOP, Object Oriented Programming) è un paradigma di programmazione che permette di definire oggetti software in grado di interagire gli uni con gli altri attraverso lo scambio di messaggi. È particolarmente adatta nei contesti in cui si possono definire delle relazioni di interdipendenza tra i concetti da modellare (contenimento, uso, specializzazione).

\letteraGlossario{P}

\definizione{PDCA}
PDCA (plan-do-check-act) è un metodo iterativo a quattro stadi usato per il controllo e il continuo miglioramento dei processi e dei prodotti.

\definizione{PDF}
Il \textbf{Portable Document Format}, comunemente abbreviato PDF, è un formato di file basato su un linguaggio di descrizione di pagina sviluppato da Adobe Systems nel 1993 per rappresentare documenti in modo indipendente dall'hardware e dal software utilizzati per generarli o per visualizzarli. 

\definizione{Package}
In alcuni linguaggi orientati agli oggetti, tra cui \glossario{Java}, è un meccanismo che permette di organizzare un insieme di classi tra loro correlate che concorrono allo stesso fine.

\definizione{Package manager}
Collezione di strumenti presenti in un sistema operativo che automatizzano il processo di installazione, aggiornamento, configurazione e rimozione dei pacchetti software in un computer.

\definizione{Parsing}
Traducibile con analisi sintattica, è il processo atto ad analizzare uno stream continuo in input in modo da determinare la sua struttura grammaticale grazie ad una data grammatica formale.

\definizione{Parser}
Strumento che si occupa di fare il \glossario{parsing}.

\definizione{Pollution}
Con tale termine ci si riferisce a tutti i file che non devono entrare nella \glossario{repository}.

\definizione{Populate}
Inserimento di dati all'interno del \glossario{database}. In questo caso all'interno dei documenti di \glossario{MongoDB}.

\definizione{Pre-commit}
Script che \glossario{verifica} i file inseriti prima di fare il commit, per evitare \glossario{pollution} dentro la \glossario{repository}.

\definizione{Proponente}
Colui che presenta una proposta, in questo caso il capitolato riguardante il progetto.
Nello specifico ci si riferisce a \Proponente{}.

\definizione{Prototipi}
Un modello approssimato o parziale del sistema che vogliamo sviluppare che simula o esegue alcune funzioni del sistema finale, realizzato allo scopo di valutarne le caratteristiche, se interno o per mostrarlo all'utente se esterno.

\definizione{Push}
Comando di \glossario{GitHub} per applicare i cambiamenti fatti nella \glossario{repository} locale alla \glossario{repository}

\definizione{Piattaforma}
In informatica il termine piattaforma indica una base software e/o hardware su cui sono sviluppate e/o eseguite applicazioni.

\letteraGlossario{Q}

\definizione{Qualità}
Insieme di caratteristiche di un'entità che ne determinano la capacità di soddisfare esigenze espresse o implicite.

\letteraGlossario{R}

\definizione{Report}
Resoconto strutturato in modo schematico.

\definizione{Repository}
Ambiente di un sistema informativo, in cui vengono gestiti i metadati, attraverso tabelle relazionali.
Nel nostro caso il sistema informativo è gestito con \glossario{GitHub}.

\definizione{Requisteak}
Software sviluppato da \GroupName{} per il tracciamento dei requisiti e relativi casi d'uso, è raggiungibile all'indirizzo \url{http://steakholders.herokuapp.com}.

\definizione{REST}
Acronimo di Representational State Transfer è un tipo di \glossario{RPC}. Si basa su un protocollo di comunicazione \glossario{stateless} di tipo \glossario{client-server}, e solitamente tale protocollo è \glossario{HTTP}.

\definizione{REST-like}
Applicazioni che non sono completamente conformi all'architettura REST.

\definizione{Registry}
Design pattern descritto in appendice della \SpecificaTecnica{}.

\definizione{RPC}
Acronimo di Remote Procedure Call, si riferisce all'attivazione di una procedura da parte di un programma, nel caso in cui tale procedura venga attivata su un computer diverso da quello sul quale il programma stesso viene eseguito.

\definizione{Ruby on Rails}
\glossario{Framework} \glossario{open source} per applicazioni web scritto in Ruby, la cui architettura è fortemente ispirata al paradigma Model-View-Controller (MVC). I suoi obiettivi sono la semplicità e la possibilità di sviluppare applicazioni di concreto interesse con meno codice rispetto ad altri \glossario{framework}. Il tutto con necessità di configurazione minimale.

\definizione{Routing}
Il routing è la funzione di un commutatore (centrale telefonica, router, switch) che decide su quale porta o interfaccia inviare un elemento di comunicazione ricevuto. Il termine è una metafora che si riferisce all'atto di instradare.

\letteraGlossario{S}

\definizione{Scaffolding}
In informatica è una procedura che automatizza la creazione di oggetti ed interfacce a partire da alcune semplici specifiche dettate dal programmatore.

\definizione{Scala}
Linguaggio di programmazione di tipo general-purpose multi-paradigma studiato per integrare le caratteristiche e funzionalità dei linguaggi orientati agli oggetti e dei linguaggi funzionali. La compilazione di codice sorgente Scala produce \glossario{Java} bytecode per l'esecuzione su una JVM.

\definizione{Scalabilità}
Capacità di un sistema di ``crescere'' o ``decrescere'' (aumentare o diminuire di scala) in funzione delle necessità e delle disponibilità.

\definizione{Schema}
Lo schema di un \glossario{database} è la struttura o organizzazione logica dei dati in esso contenuti descritta mediante l'uso di un linguaggio formale supportato da un database management system.

\definizione{SCR}
Acronimo di software change request, indica lo strumento (software) utilizzato per gestire le richieste di modifica.

\definizione{SEMAT}
Software Engineering Method and Theory è un'iniziativa nata con lo scopo di rimodellare il mondo dell'Ingegneria del Software imponendo metodi i cui fini sono la \glossario{qualità} e la disciplina.

\definizione{Server}
Componente che fornisce, attraverso una rete, un qualunque tipo di servizio o risorsa ad altri componenti tipicamente detti client.

\definizione{Show-page}
Vedi Collection-show

\definizione{Singleton}
\glossario{Design Pattern} creazionale che assicura l'esistenza di una sola istanza di una classe, e ne fornisce un punto di accesso globale.

\definizione{Slack time}
Durante la pianificazione delle attività di progetto, periodo di tempo lasciato libero da ogni attività per ammortizzare eventuali ritardi, in particolare se in presenza di forti dipendenze tra le attività.

\definizione{SMS}
Short Message Service, comunemente usato per indicare un breve messaggio di testo inviato da un telefono cellulare ad un altro.

\definizione{SMTP}
Acronimo di Simple Mail Transfer Protocol, è il protocollo standard per la trasmissione via internet di e-mail.

\definizione{Sortable}
Termine inglese traducibile con ordinabile.

\definizione{SPICE}
ISO/IEC 15504 \emph{Information technology - Process assessment}, conosciuto anche come SPICE (Software Process Improvement and Capability Determination), è un insieme di documenti di standard tecnici che fornisce informazioni generali sui concetti di valutazione dei processi e dei suoi usi nei due contesti di miglioramento dei processi e valutazione della maturità dei processi.

Riguardano i processo di sviluppo di software e le relative funzioni gestionali di azienda. In sostanza, spiega come poterne fare una valutazione.

\definizione{SQA}
Acronimo di Software Quality Assurance, insieme di azioni pro-attive che garantiscono \glossario{qualità} e rendono meno difficoltoso la \glossario{verifica}.

\definizione{Software Quality Managment Techinques}
Tecniche utili alla gestione della \glossario{qualità} mediante valutazione della qualità dei prodotti.

\definizione{Software Quality Managment}
Correlato a Software Quality Managment Techinques.

\definizione{scaffold}
Fornisce una struttura di base funzionante da cui poter procedere con incrementi.

\definizione{SCR - software change request}
Software che permette di tracciare le anomalie riscontrate nel software e nei documenti riguardanti il prodotto.

\definizione{SQL}
Acronimo per Structured Query Language, è un linguaggio standardizzato per \glossario{basi di dati} basati sul modello relazionale progettato per eseguire le funzioni \glossario{CRUD}.

\definizione{Stack}
Il termine stack o pila indica un tipo di dato astratto che viene usato in diversi contesti per riferirsi a strutture dati, le cui modalità d'accesso ai dati in essa contenuti seguono una modalità LIFO (Last In First Out), tale per cui i dati vengono estratti in ordine rigorosamente inverso rispetto a quello in cui sono stati inseriti.

\definizione{Stack tecnologico}
L'insieme delle tecnologie utilizzate dal software.

\definizione{Stakeholders}
Portatori di interesse, l'insieme di persone coinvolte nel ciclo di vita del software.

\definizione{Stand-alone}
Software in grado di funzionare da solo o in maniera indipendente da altri software, con cui può però interagire.

\definizione{Startup}
Impresa recentemente creata in cerca di un mercato. Spesso ci si riferisce ad imprese che operano nel settore informatico.

\definizione{Stateless}
Protocollo di comunicazione che tratta ogni richiesta in modo indipendente dalle richieste precedenti.

\definizione{Storage}
Si tratta di dispositivi hardware, supporti per la memorizzazione, infrastrutture e software dedicati alla memorizzazione non volatile di grandi quantità di informazioni in formato elettronico.

\definizione{Strategy}
\definizione{Design pattern} comportamentale il cui compito è isolare un algoritmo all'interno di un oggetto.

\definizione{Stub}
Porzione di codice che, dati certi input, fornisce sempre gli stessi output prestabiliti, al fine di verificare se la funzione chiamante fornisce i risultati attesi dal test.

\definizione{Sviluppatore}
Chi produce applicazioni o sistemi software.

\letteraGlossario{T}

\definizione{Tag}
Termine inglese traducibile con "etichetta". 

\definizione{Task}
Un task è un'\glossario{unità} di lavoro pianificata e assegnata. Un task può essere associato ad uno o più \glossario{ticket}, a seconda del carico di lavoro che essi richiedono.

\definizione{TeamworkPM}
Teamwork Project Manager è una \glossario{web application} che aiuta le persone coinvolte in un progetto, indipendentemente dal ruolo ricoperto a gestire e organizzare tale progetto.

\definizione{Test driver}
Clone di prova che simula la classe contenente il main. È complementare allo \glossario{Stub}.

\definizione{Three-Tier}
Si tratta di un'architettura di tipo \glossario{client-server}, in cui i processi logici, la persistenza dei dati e l'interfaccia utente sono sviluppati e mantenuti in moduli tra loro indipendenti e distribuiti. Il vantaggio che ne consegue è che un modulo può essere aggiornato o sostituito senza dover alterare gli altri.

\definizione{Ticket}
Con il termine ticket si intende un'\glossario{unità} di lavoro con cui apportare un miglioramento in un sistema. Un ticket può essere un \glossario{bug}, una richiesta di funzionalità, un compito, e così via.

\definizione{Tool}
Applicazione che svolge un determinato compito di utilità.

\definizione{Top-down}
Il nome top down significa dall'alto verso il basso: in "alto" viene posto il problema e in "basso" i sottoproblemi che lo compongono. Il nome ricorda anche una raffigurazione a piramide: l'obiettivo finale è la cima della piramide, e i sottoproblemi che lo compongono formano la base.

\letteraGlossario{U}

\definizione{UML}
Unified Modeling Language, è un linguaggio di modellazione e specifica basato sul paradigma object-oriented. 
È utilizzato per descrivere soluzioni analitiche e progettuali in modo sintetico e comprensibile a un vasto pubblico.
Ad oggi si è giunti alla versione 2.0.

\definizione{Unicode}
Sistema di codifica che assegna un numero univoco ad ogni carattere usato per la scrittura di testi, in maniera indipendente dalla lingua, dalla piattaforma informatica e dal programma utilizzato.

\definizione{Unità}
La più piccola parte di lavoro assegnabile ad un programmatore e che è utile verificare singolarmente. Un'unità raggruppa più \glossario{moduli}.

\definizione{URI}
Acronimo di Uniform Resource Identifier, è una stringa che identifica univocamente una risorsa generica che può essere un indirizzo Web, un documento, un'immagine, un file, un servizio, ecc. e la rende disponibile tramite protocolli quali HTTP, FTP, ecc.

\definizione{URL}
Acronimo di Uniform Resource Locator,  è una sequenza di caratteri che identifica univocamente l'indirizzo di una risorsa in Internet, si tratta di un termine più specifico di \glossario{URI}.

\definizione{User-friendly}
Usabilità, il grado di facilità e soddisfazione con cui si compie l'interazione tra l'uomo e uno strumento.

\definizione{Utente}
Fruitore del servizio \ProjectName{}.

\definizione{UTF-8}
Unicode Transformation Format-8 bit, è una codifica dei caratteri Unicode. 

\letteraGlossario{V}

\definizione{V8 JavaScript Engine}
Motore \glossario{JavaScript} \glossario{open source} scritto in \glossario{C++}  e sviluppato da Google, attualmente incluso in Google \glossario{Chrome}.

\definizione{Validazione}
Controllo effettuato sul software, per controllare se tutti i requisiti previsti sono stati coperti.

\definizione{Verifica}
Ricerca la consistenza, correttezza e completezza.

\definizione{Viste}
Interfaccia utente nel modello \glossario{MVC}.

\letteraGlossario{W}

\definizione{W3C}
Abbreviazione per World Wide Web Consortium, è un'organizzazione non governativa internazionale che ha come scopo quello di sviluppare e standardizzare tutte le potenzialità del World Wide Web.

\definizione{Web application}
Applicazione accessibile via web, che offre determinati servizi all'utente.

\letteraGlossario{X}

\definizione{XML}
Sigla di eXtensible Markup Language, è un linguaggio di \glossario{markup}.
%\appendix
%\include{appendice-A.tex}

\end{document}
